\documentclass[11pt,a4paper]{article}

\usepackage[utf8]{inputenc}
\usepackage[T1]{fontenc}
\usepackage{lmodern}
\usepackage{textcomp}
\usepackage[margin=1in]{geometry}
\usepackage{amsmath}
\usepackage{amssymb}
\usepackage{booktabs}
\usepackage{longtable}
\usepackage{array}
\usepackage{xcolor}
\usepackage{hyperref}

\hypersetup{
  colorlinks=true,
  linkcolor=black,
  citecolor=black,
  urlcolor=blue
}

\setlength{\parindent}{0pt}
\setlength{\parskip}{0.72em}
\renewcommand{\arraystretch}{1.18}

\newcommand{\code}[1]{\texttt{#1}}
\newcommand{\claimnote}[1]{%
  \vspace{0.5em}
  \noindent\fbox{%
    \begin{minipage}{0.94\linewidth}
    #1
    \end{minipage}}
  \vspace{0.5em}
}

\title{\textbf{DAD FieldWorks}\\
{\large A Reusable Evidence Contract Architecture for Validation Aware Computational Electromagnetics Software}}
\author{Harun Aktas}
\date{Technical Whitepaper / Public Specification Draft\\Version 0.5 draft\\May 2026}

\begin{document}

\maketitle

\begin{center}
\textbf{Claim status:} No external validation claim. No production readiness claim.\\
Copyright © 2026 Harun Aktas. All rights reserved.\\
No license is granted for reuse, redistribution, or derivative work without written permission from the copyright holder.
\end{center}

\newpage
\begin{abstract}
DAD FieldWorks is a private research and development project for evidence-aware computational electromagnetics software. This public specification draft describes a reusable evidence contract architecture that keeps numerical execution, evidence records, gate rows, decision classes, failure classes, reference comparisons, risks, required future work, and claim boundaries in one structured result model.

The document summarizes public-safe current status and selected internal evidence-path results. It does not publish private source code, private tests, notebooks, internal project-control files, or private artifact files. Numerical tables in this draft are internal research evidence and diagnostic comparisons only. They are not external validation evidence, not production readiness claims, and not commercial solver equivalence claims.
\end{abstract}

\tableofcontents
\newpage

\section{Introduction}
\label{sec:introduction}

Computational electromagnetics software often produces numerical output long before the output is safe to use as a claim. A field plot, residual, frequency estimate, or comparison table can be technically useful while still being insufficient for release, certification, external agreement, or engineering reliance. DAD FieldWorks treats this gap as an architectural problem.

The central idea in this whitepaper is an evidence contract: every solver-facing result should carry not only values, but also structured evidence, gate rows, decision labels, failure labels, reference-comparison rows, risks, required future work, and explicit claim boundaries. This keeps numerical success from silently becoming a claim that the available evidence does not support.

Version 0.5 extends the public whitepaper source with current public-safe project status, a C++ evidence contract foundation, current PEC cavity diagnostic evidence, and provenance for the public solver generated hero visualization.

\claimnote{\textbf{Public boundary.} This document is a public architecture and status draft. It is not validation evidence, not a production readiness claim, not a product release, and not an external solver equivalence claim.}

\section{Motivation}
\label{sec:motivation}

Electromagnetic solver work contains several distinct layers:

\begin{itemize}
  \item numerical implementation,
  \item problem definition and geometry metadata,
  \item grid, mesh, basis, material, and boundary conventions,
  \item source and observable definitions,
  \item analytical or alternative computational code comparison rows,
  \item residual and convergence diagnostics,
  \item gate decisions and failure classification,
  \item risk and required future work,
  \item claim boundary review.
\end{itemize}

These layers should not be compressed into a single Boolean. A solver can execute without being comparison-ready. A comparison can be internally useful without allowing an external validation claim. A clean residual can coexist with a boundary convention ambiguity. A visualization can be solver generated and still remain diagnostic only.

The evidence contract architecture is motivated by that separation. It encourages each solver family to return structured evidence that makes the current decision state reviewable, repeatable, and conservative.

\section{Relationship to CEM Verification and Validation Practice}
\label{sec:relationship-to-vvuq}

Classical computational electromagnetics verification and validation practice distinguishes implementation correctness, numerical convergence, reference comparison, uncertainty, and model adequacy. The DAD FieldWorks evidence contract architecture does not replace those practices. It provides a reusable result shape for carrying their status through solver workflows.

The architecture is designed to support:

\begin{itemize}
  \item analytical reference checks for simple cases,
  \item internal code-verification diagnostics,
  \item grid, mesh, and basis refinement evidence,
  \item source, probe, and observable metadata,
  \item bounded artifact reviews,
  \item alternative computational code comparison workflows,
  \item explicit limitations and claim boundaries.
\end{itemize}

The architecture deliberately separates the phrase ``computed successfully'' from any phrase implying release status. Computed values must be interpreted through the evidence rows and gate rows that produced them.

\section{Evidence Contract Architecture}
\label{sec:evidence-contract-architecture}

An evidence contract is a structured result object with four responsibilities:

\begin{enumerate}
  \item carry solver and case metadata,
  \item carry numerical values and evidence rows,
  \item carry gate decisions and failure classes,
  \item carry claim boundaries and future work rows.
\end{enumerate}

The core public-safe type vocabulary used in the current architecture is:

\begin{longtable}{p{0.30\linewidth}p{0.62\linewidth}}
\toprule
\textbf{Record type} & \textbf{Purpose}\\
\midrule
\endhead
\code{EvidenceRow} & General structured evidence record for values, metrics, provenance, assumptions, or diagnostic facts.\\
\code{GateRow} & Named decision gate with pass flag, value, threshold, or notes.\\
\code{DecisionClass} & Controlled decision label for the result state.\\
\code{FailureClass} & Controlled failure or blocked-state label.\\
\code{ClaimBoundary} & Explicit allowed and disallowed public claim row.\\
\code{ReferenceComparisonRow} & Analytical, manufactured, or alternative computational code comparison row.\\
\code{RiskRow} & Technical risk, mitigation, and residual uncertainty row.\\
\code{RequiredFutureWorkRow} & Work required before a stronger claim or next task can be opened.\\
\code{SolverMetadata} & Solver family, solver name, method, version, and case identity.\\
\code{ValidationArtifactMetadata} & Identifiers for reviewable artifacts or evidence bundles.\\
\code{ResultContractSummary} & Summary decision, failure class, and conservative claim flags.\\
\code{SolverEvidenceResult} & Top-level result container collecting metadata, rows, gates, comparisons, risks, future work, and claim boundaries.\\
\bottomrule
\end{longtable}

The model is intentionally reusable. It can represent a scalar eigenmode prototype, a time-domain FDTD diagnostic, a FEM Helmholtz check, a Method of Moments antenna workflow, or an external comparison workflow without changing the result contract shape.

\section{Evidence Row Model}
\label{sec:evidence-row-model}

An \code{EvidenceRow} records a specific fact that may be useful for review. Typical fields include an identifier, evidence class, quantity name, value, units, method note, and interpretation note.

Examples of evidence row content include:

\begin{itemize}
  \item grid dimensions and physical dimensions,
  \item material and boundary convention records,
  \item eigenfrequency estimates and residual norms,
  \item FDTD source, probe, and ringdown metadata,
  \item field component availability rows,
  \item artifact readability rows,
  \item execution-boundedness rows.
\end{itemize}

Evidence rows are not themselves claims. They become actionable only when combined with gate rows, risk rows, and claim boundaries.

\section{Gate Row Model}
\label{sec:gate-row-model}

A \code{GateRow} is a named check. It should be specific enough that a reviewer can see why a decision passed, failed, or remained blocked.

\begin{longtable}{p{0.28\linewidth}p{0.24\linewidth}p{0.38\linewidth}}
\toprule
\textbf{Gate family} & \textbf{Typical state} & \textbf{Example interpretation}\\
\midrule
\endhead
Execution gate & pass/fail & Did the bounded computation execute without non-finite output?\\
Residual gate & pass/fail & Is the residual below the stated internal threshold?\\
Reference gate & pass/fail/blocked & Is a reference comparison present and scoped?\\
Metadata gate & pass/fail/blocked & Are grid, boundary, source, and observable conventions recorded?\\
Claim gate & pass/fail & Are external validation and production flags still conservative?\\
Artifact gate & pass/fail & Are review artifacts readable and bounded?\\
\bottomrule
\end{longtable}

The gate row model is deliberately granular. A result can pass residual gates and fail metadata gates. That is a useful state, not an error in the architecture.

\section{Decision and Failure Classification}
\label{sec:decision-and-failure-classification}

\code{DecisionClass} and \code{FailureClass} make result state explicit. They should not be anonymous strings hidden in log output. They should be part of the returned result contract.

Examples of conservative decision classes include:

\begin{itemize}
  \item scalar grid refinement converging,
  \item same-grid scalar interpretation inconclusive,
  \item records-only policy audit passed,
  \item solver health issues found,
  \item comparison blocked pending metadata closure.
\end{itemize}

Examples of conservative failure classes include:

\begin{itemize}
  \item none,
  \item source required,
  \item solver health issues found,
  \item validation data required,
  \item mapping ambiguity unresolved,
  \item comparison blocked.
\end{itemize}

The project currently treats decision and failure labels as typed string wrappers in the C++ evidence contract foundation. This allows solver-specific vocabularies while preventing decision state from dissolving into unstructured notes.

\section{Reference Comparison Model}
\label{sec:reference-comparison-model}

A \code{ReferenceComparisonRow} records comparison context. It should not be only a numerical difference. A useful comparison row should include:

\begin{itemize}
  \item reference type,
  \item reference quantity,
  \item computed quantity,
  \item absolute and relative difference where meaningful,
  \item tolerance policy,
  \item scope of applicability,
  \item claim boundary.
\end{itemize}

Reference comparisons may involve analytical formulas, manufactured solutions, regression baselines, or alternative computational code comparisons. Each comparison must keep its scope. A comparison for one mode, mesh, observable, or source does not automatically generalize to a solver family.

\section{Claim Boundary Model}
\label{sec:claim-boundary-model}

The claim boundary model is where the architecture prevents accidental overstatement.

\begin{longtable}{p{0.34\linewidth}p{0.56\linewidth}}
\toprule
\textbf{Boundary} & \textbf{Current conservative rule}\\
\midrule
\endhead
External validation & Numerical success alone does not create an external validation claim.\\
Production readiness & Execution, residuals, and internal comparison rows do not create a production readiness claim.\\
Commercial solver equivalence & A scoped comparison row does not create a general commercial solver equivalence claim.\\
Measurement comparison & No measurement validation claim exists without traceable measurement artifacts and gates.\\
Scalar-to-vector transfer & Scalar Helmholtz prototypes do not establish a full vector Maxwell eigenmode solver claim.\\
Visualization & Solver generated visualization is diagnostic presentation only.\\
\bottomrule
\end{longtable}

The current conservative defaults are:

\begin{center}
\begin{tabular}{ll}
\toprule
\textbf{Flag} & \textbf{Default}\\
\midrule
\code{ExternallyValidatedQ} & \code{False}\\
\code{ProductionAllowedQ} & \code{False}\\
\bottomrule
\end{tabular}
\end{center}

These flags do not become true by numerical success alone. They require separate evidence, acceptance gates, and claim review.

\section{Reusable VVUQ Patterns}
\label{sec:reusable-vvuq-patterns}

The current evidence chain suggests reusable VVUQ patterns:

\begin{itemize}
  \item \textbf{Bounded execution:} limit problem size, stored fields, time histories, and artifacts.
  \item \textbf{Records-only planning:} plan comparison or implementation gates before running new numerics.
  \item \textbf{Source and observable discipline:} record source type, source position, probe position, component, window, and observable.
  \item \textbf{Reference separation:} distinguish analytical reference, internal diagnostic estimate, and external comparison workflow.
  \item \textbf{Ambiguity rows:} preserve unresolved mapping, component, source, probe, and boundary ambiguity.
  \item \textbf{No silent promotion:} do not promote an internal frequency candidate into a claim without explicit gates.
  \item \textbf{Claim defaults:} keep external validation and production flags false unless a separate claim review changes them.
\end{itemize}

These patterns are reusable because they are independent of one solver method.

\section{Application to Solver Families}
\label{sec:application-to-solver-families}

\subsection{Eigenmode Solvers}

Eigenmode evidence should carry geometry, material, boundary, basis or grid convention, eigenvalue formulation, residual definition, normalization policy, nullspace treatment, and reference comparison rows. Scalar Helmholtz prototypes can be useful stepping stones, but they must not be used as proof of a full vector Maxwell eigenmode chain.

\subsection{FDTD Solvers}

FDTD evidence should carry time step, Courant factor, grid spacing, source type, source location, probe location, boundary treatment, component family, stored observable, ringdown window, peak-selection policy, and field snapshot limitations. Internal PEC cavity FDTD evidence remains diagnostic and comparison-blocked where mode ambiguity, component family disagreement, or metadata closure is unresolved.

\subsection{FEM Helmholtz Solvers}

FEM Helmholtz evidence should carry weak form, element family, mesh, material map, boundary condition, residual norm, manufactured-solution or analytical reference rows, and mesh refinement gates. A manufactured-solution check remains internal verification unless claim rows support more.

\subsection{Method of Moments Antenna Solvers}

Method of Moments evidence should carry integral equation choice, basis functions, mesh or segment metadata, feed model, impedance or radiation observable, convergence rows, and reference comparison rows. Feed and observable definitions are as important as numerical convergence.

\subsection{External Comparison Workflows}

External comparison workflows should carry artifact metadata, same-problem mapping evidence, normalization gates, quantity definitions, tolerance policy, risk rows, and claim boundary rows. The comparison should be scoped to the case, observable, and artifact quality available.

\section{C++ Evidence Contract Foundation}
\label{sec:cpp-evidence-contract-foundation}

The current C++ direction is a lightweight evidence-contract foundation rather than a released solver core. Public-safe current state:

\begin{itemize}
  \item solver-agnostic C++ value types exist for the evidence contract vocabulary,
  \item the contract centers on \code{SolverEvidenceResult},
  \item \code{ResultContractSummary} keeps conservative default claim flags,
  \item local C++ contract tests exercise the generic contract layer,
  \item a narrow non-numerical scalar Helmholtz demo result wrapper demonstrates how a future wrapper can populate metadata, evidence rows, gates, comparison readiness rows, risks, claim boundaries, and required future work,
  \item no C++ solver core is released in this public repository,
  \item no C++ eigenmode, FDTD, FEM, or Method of Moments numerical kernel is claimed as public release material.
\end{itemize}

The current C++ foundation is important because it moves evidence contracts toward future solver wrappers without requiring private implementation details in public material. It is infrastructure and shape discipline, not solver validation evidence.

\section{Current Public Technical Status}
\label{sec:current-public-technical-status}

DAD FieldWorks remains a private research and development project. The public showcase contains selected architecture material, public status notes, a public whitepaper PDF, public companion notes, and solver generated diagnostic visualization assets.

Current public-safe status from the private project state can be summarized as follows:

\begin{itemize}
  \item Evidence contract architecture is the central reusable pattern.
  \item Conservative claim defaults remain active: \code{ExternallyValidatedQ = False} and \code{ProductionAllowedQ = False}.
  \item Scalar PEC cavity Helmholtz Route A refinement evidence is internally convergent for the bounded scalar model.
  \item Same-grid scalar PEC cavity diagnostics remain inconclusive because mapping policy ambiguity is material.
  \item FDTD PEC cavity internal evidence contains useful component and frequency-family diagnostics, but mode ambiguity and comparison-readiness gaps remain.
  \item Records-only Route A/C policy closure for PEC cavity native-Yee planning is complete at the policy layer, while derivative boundary policy, normal component disposition, residual policy, nullspace policy, mass/generalized eigenproblem policy, and native curl-curl implementation remain blocked.
  \item C++ currently provides evidence-contract infrastructure and a non-numerical demo wrapper, not a released solver core.
\end{itemize}

\section{Case Study: PEC Cavity Evidence Chain}
\label{sec:pec-cavity-evidence-chain}

The PEC cavity case is useful because it is simple enough to have an analytical reference while still exposing real evidence-contract issues: grid interpretation, field component conventions, boundary policy, source/probe dependence, frequency-family ambiguity, and claim boundaries.

\subsection{Reference Case}

The current public-safe reference case is a rectangular PEC cavity with:

\begin{center}
\begin{tabular}{ll}
\toprule
\textbf{Quantity} & \textbf{Value}\\
\midrule
\(a\) & \(0.08\ \mathrm{m}\)\\
\(b\) & \(0.084\ \mathrm{m}\)\\
\(d\) & \(0.084\ \mathrm{m}\)\\
\((m,n,p)\) & \((1,1,1)\)\\
\(c_0\) & \(299792458\ \mathrm{m/s}\)\\
\text{Analytical reference frequency} & \(3.143165987503105\ \mathrm{GHz}\)\\
\bottomrule
\end{tabular}
\end{center}

For a rectangular cavity reference, the frequency form is:

\[
f_{mnp} = \frac{c_0}{2}\sqrt{\left(\frac{m}{a}\right)^2+\left(\frac{n}{b}\right)^2+\left(\frac{p}{d}\right)^2}.
\]

This analytical reference is a scoped reference for the case. It does not by itself validate any solver family.

\subsection{Scalar Helmholtz Route A Grid Refinement}

The scalar Helmholtz Route A refinement evidence is a bounded homogeneous-Dirichlet scalar finite-difference diagnostic. It is useful scalar convergence evidence only.

\begin{longtable}{ccccc}
\toprule
\textbf{Level} & \textbf{Grid} & \textbf{Frequency (GHz)} & \textbf{Relative error} & \textbf{Residual norm}\\
\midrule
\endhead
0 & \(\{9,8,8\}\) & 3.128307359058 & 0.0047272809 & \(2.63\times 10^{-15}\)\\
1 & \(\{18,16,16\}\) & 3.139012096545 & 0.0013215627 & \(1.29\times 10^{-14}\)\\
2 & \(\{27,24,24\}\) & 3.141247276062 & 0.0006104391 & \(1.68\times 10^{-14}\)\\
\bottomrule
\end{longtable}

The relative error decreases monotonically for these three bounded levels. The optional larger row remains outside the bounded execution status used for this public draft.

\subsection{Same-Grid Scalar Eigenmode Diagnostic}

The same-grid scalar diagnostic compares multiple interpretations of grid metadata and scalar boundary convention. The numerical eigenvalue solves are clean, but interpretation remains the limiting factor.

\begin{longtable}{ccccc}
\toprule
\textbf{Policy} & \textbf{Matrix size} & \textbf{Frequency (GHz)} & \textbf{Relative error} & \textbf{Residual norm}\\
\midrule
\endhead
1 & \(504\times 504\) & 2.778104015200 & 0.1161446687 & \(1.83\times 10^{-15}\)\\
2 & \(336\times 336\) & 3.122528486251 & 0.0065658325 & \(1.14\times 10^{-15}\)\\
3 & \(504\times 504\) & 3.126946915965 & 0.0051601066 & \(2.37\times 10^{-15}\)\\
\bottomrule
\end{longtable}

The same-grid scalar diagnostic remains inconclusive because mapping policy ambiguity is material. Policy 1 is clean numerically but has a large relative error. Policies 2 and 3 are plausible, with Policy 3 closest to the analytical reference in this pass. No mapping policy is promoted by this evidence alone.

\subsection{Internal Frequency Comparison}

\begin{longtable}{p{0.46\linewidth}cc}
\toprule
\textbf{Route or diagnostic} & \textbf{Frequency (GHz)} & \textbf{Relative error}\\
\midrule
\endhead
Analytical reference & 3.143165987503 & 0\\
Route A scalar Helmholtz Level 0 & 3.128307359058 & 0.004727281\\
Route A scalar Helmholtz Level 1 & 3.139012096545 & 0.001321563\\
Route A scalar Helmholtz Level 2 & 3.141247276062 & 0.000610439\\
Same-grid scalar Policy 2 & 3.122528486251 & 0.006565833\\
Same-grid scalar Policy 3 & 3.126946915965 & 0.005160107\\
Numerical dispersion diagnostic & 3.126477519038 & 0.005309445\\
FDTD EzHxFamily & 2.950641806786 & 0.061251675\\
FDTD ExFamily & 3.631559146814 & 0.155382554\\
\bottomrule
\end{longtable}

This is an internal evidence path comparison. It is not external validation evidence. The FDTD family rows remain diagnostic because component-family disagreement, source/probe effects, peak-selection policy, and mode identification remain claim boundaries.

\subsection{Current PEC Cavity Status}

Current public-safe status for the PEC cavity evidence chain:

\begin{itemize}
  \item Route A/C closure is records-closed as a policy layer.
  \item Normal electric component disposition remains unresolved.
  \item Boundary derivative policy remains partial.
  \item Residual metric policy remains partial.
  \item Mass or generalized eigenproblem policy remains partial.
  \item Nullspace and spurious-mode policy remains partial.
  \item Native Yee curl-curl eigenmode prototype planning remains blocked.
  \item CPML support is not claimed for the normal solver path.
  \item No internal frequency candidate is promoted to a validation claim.
\end{itemize}

\section{Solver Generated Public Diagnostic Visualization}
\label{sec:solver-generated-public-diagnostic-visualization}

The public showcase includes a solver generated diagnostic visualization. Its public-safe provenance is:

\begin{longtable}{p{0.34\linewidth}p{0.56\linewidth}}
\toprule
\textbf{Item} & \textbf{Public-safe description}\\
\midrule
\endhead
Field family & 2D FDTD TMz\\
Updated fields & Ez, Hx, Hy\\
Rendered quantity & Signed Ez\\
Object & Rectangular PEC-like object from the simulation mask\\
PEC condition & Ez held at zero inside the PEC mask\\
Source & Gaussian-modulated sine soft point source\\
Boundary handling & Graded loss layer with simple absorbing edge copy\\
External images & none\\
Screenshots & none\\
AI image generation & none\\
Private source code in public repo & none\\
Claim boundary & diagnostic visualization only; not validation evidence; not production readiness\\
\bottomrule
\end{longtable}

The visualization is intended to communicate field diagnostic character: incoming wave, PEC interaction, reflection, scattering, shadowing, and time evolution. It does not establish validation, comparison agreement, or release readiness.

\section{Advantages of the Proposed Architecture}
\label{sec:advantages}

The proposed architecture has several practical advantages:

\begin{itemize}
  \item \textbf{Reviewability:} decisions and limitations are inspectable in structured rows.
  \item \textbf{Reusability:} one contract shape can serve multiple solver families.
  \item \textbf{Conservatism:} claim flags default to false and require explicit evidence to change.
  \item \textbf{Separation of concerns:} numerical execution, reference comparison, and public claims stay distinct.
  \item \textbf{Migration path:} mature research evidence patterns can be carried into future C++ wrappers.
  \item \textbf{Failure clarity:} blocked states become useful information rather than hidden weakness.
\end{itemize}

\section{Limitations and Claim Boundaries}
\label{sec:limitations}

This draft makes the following limitations explicit:

\begin{itemize}
  \item no external validation claim,
  \item no production readiness claim,
  \item no commercial solver equivalence claim,
  \item no formal external solver agreement or certification claim,
  \item no measurement validation claim,
  \item no full-vector Maxwell eigenmode solver claim from scalar prototypes,
  \item no native Yee curl-curl solver claim from the current scalar prototypes,
  \item no external FDTD validation claim,
  \item same-grid scalar mapping ambiguity remains open,
  \item FDTD PEC cavity component-family ambiguity remains open,
  \item public visualization is diagnostic only,
  \item C++ evidence contracts are infrastructure and do not release a numerical solver core.
\end{itemize}

\section{Roadmap}
\label{sec:roadmap}

\subsection{Near Term}

\begin{itemize}
  \item keep validation and production flags conservative,
  \item continue same-grid scalar interpretation review,
  \item document grid convention and boundary mapping requirements,
  \item close boundary derivative route policy at records level before implementation,
  \item continue C++ compilation and contract testing for the evidence foundation,
  \item bind only narrow, reviewable solver wrappers to \code{SolverEvidenceResult}.
\end{itemize}

\subsection{Mid Term}

\begin{itemize}
  \item define native Yee field component offsets and boundary masks,
  \item document native curl and curl-curl metadata requirements,
  \item plan a native Yee curl-curl eigenmode prototype only after metadata closure,
  \item develop independent FDTD reference cases that do not depend on PEC feature tracking alone,
  \item improve source, probe, ringdown, and peak-selection evidence rows,
  \item define FEM Helmholtz and Method of Moments evidence contract wrappers.
\end{itemize}

\subsection{Long Term}

\begin{itemize}
  \item mature a full eigenmode solver evidence chain,
  \item establish independent FDTD validation cases under traceable gates,
  \item build external comparison pipelines with artifact metadata and same-problem mapping gates,
  \item migrate mature research kernels into C++ only when evidence contracts are mandatory,
  \item maintain claim boundary review as a release gate for each solver family.
\end{itemize}

\section{Conclusion}
\label{sec:conclusion}

DAD FieldWorks uses evidence contracts to keep computational electromagnetics research honest about what has been computed, what has been checked, what remains ambiguous, and what cannot yet be claimed. The architecture is reusable across eigenmode, FDTD, FEM Helmholtz, Method of Moments, and external comparison workflows.

The current public v0.5 draft documents meaningful progress: a reusable evidence row and gate row vocabulary, conservative claim defaults, scalar PEC cavity refinement evidence, same-grid scalar ambiguity, FDTD PEC cavity diagnostic family comparisons, a public solver generated visualization, and a C++ evidence contract foundation. The limitations are equally important: no external validation claim, no production readiness claim, no commercial solver equivalence claim, no private source code publication, and no solver-core release in the public showcase.

\section{References}
\label{sec:references}

\begin{thebibliography}{9}

\bibitem{yee1966}
K. S. Yee,
``Numerical solution of initial boundary value problems involving Maxwell's equations in isotropic media,''
\emph{IEEE Transactions on Antennas and Propagation}, vol. 14, no. 3, pp. 302--307, 1966.

\bibitem{taflove}
A. Taflove and S. C. Hagness,
\emph{Computational Electrodynamics: The Finite-Difference Time-Domain Method}.
Artech House, third edition, 2005.

\bibitem{sullivan}
D. M. Sullivan,
\emph{Electromagnetic Simulation Using the FDTD Method}.
IEEE Press, 2000.

\bibitem{bondeson}
A. Bondeson, T. Rylander, and P. Ingelstrom,
\emph{Computational Electromagnetics}.
Springer, 2005.

\bibitem{balanis}
C. A. Balanis,
\emph{Advanced Engineering Electromagnetics}.
Wiley, second edition, 2012.

\bibitem{oberkampf}
W. L. Oberkampf and C. J. Roy,
\emph{Verification and Validation in Scientific Computing}.
Cambridge University Press, 2010.

\bibitem{roache}
P. J. Roache,
\emph{Verification and Validation in Computational Science and Engineering}.
Hermosa Publishers, 1998.

\end{thebibliography}

\end{document}
